\chapter*{Conclusione} %l'asterisco dopo chapter serve per visualizzare il capitolo come "non numerato"

In questa tesi è stato analizzato il problema della rappresentazione di una struttura dati gerarchica classica, nello specifico un max-heap, all'interno del modello di computazione quantistica. L'obiettivo principale del lavoro è stato studiare quali componenti siano necessarie per definire una struttura equivalente nel paradigma quantistico e verificare se la proprietà fondamentale di ordinamento possa essere espressa attraverso operatori e circuiti quantistici.

La prima parte del lavoro ha evidenziato come una trasposizione diretta delle strutture dati classiche non sia immediatamente possibile. In particolare, la proprietà di max-heap richiede una relazione d'ordine tra gli elementi contenuti nella struttura, mentre gli stati quantistici arbitrari non possiedono una relazione di ordinamento naturale. Per questo motivo, il modello adottato mantiene una rappresentazione classica dei valori memorizzati, introducendo invece il paradigma quantistico nella gestione degli indirizzi e nelle operazioni di accesso e verifica.

Successivamente sono state definite le primitive necessarie alla costruzione della struttura proposta. In particolare, sono stati implementati due differenti modelli di accesso alla memoria, uno basato su una codifica circuitale assimilabile ad una Quantum Read Only Memory e uno basato su rotazioni controllate del registro dei dati. Entrambe le soluzioni consentono di realizzare l'operazione di accesso indicizzato alla memoria quantistica, mostrando tuttavia differenti compromessi in termini di semplicità implementativa e costo circuitale.

Sono stati inoltre introdotti gli operatori necessari alla navigazione della struttura, permettendo il calcolo degli indirizzi del nodo padre e dei figli a partire dalla rappresentazione binaria dell'indice. Su queste primitive è stato costruito un verificatore della proprietà di max-heap, basato sul confronto tra ciascun nodo e i suoi figli diretti. La verifica globale della struttura è stata successivamente formulata come un problema di ricerca di indirizzi che violano tale proprietà, consentendo l'applicazione dell'algoritmo di Grover.

I risultati sperimentali ottenuti mediante simulazione hanno confermato il corretto funzionamento delle primitive proposte. In particolare, gli esperimenti condotti su istanze di heap di dimensione 7 elementi hanno mostrato il corretto comportamento dei moduli di accesso alla memoria, del verificatore della proprietà e dell'oracolo utilizzato nella ricerca quantistica, sia in presenza di strutture valide sia in presenza di uno o più elementi che violano il vincolo di max-heap.

Il lavoro presenta tuttavia alcuni limiti. La struttura proposta non implementa ancora tutte le operazioni tipiche di un heap classico, come inserimento, estrazione del massimo o procedura di max-heapify. Inoltre, le valutazioni effettuate sono basate su simulazioni di circuiti quantistici e non considerano gli effetti del rumore e delle limitazioni hardware dei dispositivi quantistici attuali. Infine, il vantaggio teorico ottenuto tramite Grover riguarda il numero di interrogazioni dell'oracolo e non implica necessariamente un vantaggio computazionale complessivo, il quale dipende dal costo della realizzazione della memoria quantistica e del verificatore.

Come sviluppi futuri, risulta interessante lo studio di una procedura di max-heapify quantistica, che permetterebbe di estendere la struttura dalla sola verifica dell'invariante alla sua manutenzione dinamica. Ulteriori direzioni di ricerca riguardano l'ottimizzazione delle primitive di accesso e l'indagine di strutture dati gerarchiche più complesse nel modello quantistico.